\chapter*{Abstract}
\addcontentsline{toc}{chapter}{Abstract}

This report presents the design and implementation of Qatra, a distributed platform intended to improve blood-donation coordination and emergency response. The solution connects donors, donation-center staff, center administrators, and platform administrators through unified workflows for registration, health profiling, appointment scheduling, eligibility follow-up, emergency creation, donor matching, and real-time notification. The platform follows a microservices-oriented architecture centered on two Java 21 and Spring Boot services, an Angular 20 web application, PostgreSQL persistence, Kafka event exchange, Redis, and a containerized operational environment. Kubernetes manifests, Nginx routing, GitHub Actions, Prometheus, Grafana, Loki, Promtail, and Jaeger support deployment and observability concerns. Kanban was used to organize incremental delivery and maintain traceability between requirements, implementation, and validation. The report emphasizes the domain architecture, emergency matching logic, event-driven notification pipeline, security model, and the limits of the available verification evidence. Qatra demonstrates how modern distributed-system practices can support a socially significant domain while preserving maintainability, auditability, and operational visibility.

\textbf{Keywords:} blood donation, emergency response, microservices, Spring Boot, Angular, Kafka, Kubernetes, Kanban, observability.

\chapter*{Résumé}
\addcontentsline{toc}{chapter}{Résumé}

Ce rapport présente la conception et la réalisation de Qatra, une plateforme distribuée destinée à améliorer la coordination des dons de sang et la réponse aux urgences. La solution relie les donneurs, le personnel des centres, les administrateurs de centres et les administrateurs de la plateforme à travers des parcours unifiés d'inscription, de profil médical, de prise de rendez-vous, de suivi d'éligibilité, de création d'urgence, de mise en correspondance des donneurs et de notification en temps réel. La plateforme adopte une architecture orientée microservices comprenant deux services Java 21 et Spring Boot, une application web Angular 20, PostgreSQL, Kafka, Redis et un environnement opérationnel conteneurisé. Des manifestes Kubernetes, Nginx, GitHub Actions, Prometheus, Grafana, Loki, Promtail et Jaeger couvrent le déploiement et l'observabilité. La méthode Kanban structure la livraison incrémentale et la traçabilité entre exigences, réalisation et validation. Le rapport met l'accent sur l'architecture métier, le moteur de mise en correspondance d'urgence, les notifications événementielles, la sécurité et les limites des preuves de validation disponibles.

\textbf{Mots-clés :} don de sang, réponse d'urgence, microservices, Spring Boot, Angular, Kafka, Kubernetes, Kanban, observabilité.

